% fig_dataset_examples.tex -- dataset examples: one schematic frame-strip
% per system, showing what a conditioning sequence + swept theta produces.
% Draft only. Real grids/units taken from main.tex's Table (tab:suite).
%
% Each row: a short strip of 4 schematic frames (ball/disk on plain
% background against ground/pivot/wall as appropriate) at increasing time,
% with theta and its true grid values annotated. This is meant to show a
% reader concretely what "the PhysWeep suite" renders, since the paper's
% own suite table (tab:suite) is text-only.
%
% REQUIRES: \usetikzlibrary{decorations.pathmorphing} (for the spring
% coil in the spring-mass row) -- not loaded by main.tex's preamble.tex,
% add it if copying this figure in.

\begin{tikzpicture}[
  font=\small,
  framebox/.style={draw=black!40, fill=black!3, minimum width=1.55cm,
                    minimum height=1.55cm, line width=0.6pt},
  disk/.style={fill=RoyalBlue!75!black, draw=none},
  ground/.style={black!55, line width=1.0pt},
  rowlbl/.style={font=\bfseries\small, anchor=east},
  gridlbl/.style={font=\scriptsize, black!55, anchor=west, align=left}
]

% column x-centers for the 4-frame strip
\def\colx{0,1.75,3.5,5.25}
\def\rowlabelx{-2.6}
\def\stripstartx{-1.0}

% ============ ROW 1: projectile ============
\node[rowlbl] at (\rowlabelx,0) {Projectile};
\foreach \i/\x in {1/0,2/1.75,3/3.5,4/5.25}{
  \node[framebox] (p\i) at (\x,0) {};
}
\draw[ground] ([yshift=-0.7cm]p1.south west) -- ([yshift=-0.7cm]p4.south east);
% falling arc across 4 frames, y decreasing then bottoming out
\fill[disk] (p1.center)+(0,0.55) circle (2.6pt);
\fill[disk] (p2.center)+(0,0.15) circle (2.6pt);
\fill[disk] (p3.center)+(0,-0.35) circle (2.6pt);
\fill[disk] (p4.center)+(0,-0.6) circle (2.6pt);
\node[gridlbl] at (6.3,0) {$\theta{=}g$: in $\{5,7,9.8,12,15\}$\\out $\{1.6,3,20,25\}$ m/s$^2$};

% ============ ROW 2: damped pendulum ============
\node[rowlbl] at (\rowlabelx,-2.0) {Damped pendulum};
\foreach \i/\x in {1/0,2/1.75,3/3.5,4/5.25}{
  \node[framebox] (d\i) at (\x,-2.0) {};
}
\foreach \i/\x/\ang in {1/0/35,2/1.75/15,3/3.5/-20,4/5.25/-8}{
  \fill[black!55] ([yshift=0.72cm]d\i.center) circle (1.0pt);
  \draw[black!65,line width=0.9pt] ([yshift=0.72cm]d\i.center) -- ++({-90+\ang}:0.62cm);
  \fill[disk] ([yshift=0.72cm]d\i.center) ++({-90+\ang}:0.62cm) circle (2.4pt);
}
\node[gridlbl] at (6.3,-2.0) {$\theta{=}\omega$: in $\{2,3,4,5\}$\\out $\{0.5,1,8,12\}$ rad/s\\(2nd axis: $\zeta$, envelope decay)};

% ============ ROW 3: bouncing ball ============
\node[rowlbl] at (\rowlabelx,-4.0) {Bouncing ball};
\foreach \i/\x in {1/0,2/1.75,3/3.5,4/5.25}{
  \node[framebox] (b\i) at (\x,-4.0) {};
}
\draw[ground] ([yshift=-0.7cm]b1.south west) -- ([yshift=-0.7cm]b4.south east);
\fill[disk] (b1.center)+(0,0.6) circle (2.6pt);
\fill[disk] (b2.center)+(0,-0.68) circle (2.6pt);
\fill[disk] (b3.center)+(0,0.25) circle (2.6pt);
\fill[disk] (b4.center)+(0,-0.68) circle (2.6pt);
\node[gridlbl] at (6.3,-4.0) {$\theta{=}e$: in $\{.6,.7,.8,.9\}$\\out $\{.2,.35,.97\}$ (dimensionless)};

% ============ ROW 4: spring-mass ============
\node[rowlbl] at (\rowlabelx,-6.0) {Spring--mass};
\foreach \i/\x in {1/0,2/1.75,3/3.5,4/5.25}{
  \node[framebox] (s\i) at (\x,-6.0) {};
}
\foreach \i/\x/\yy in {1/0/0.55,2/1.75/-0.15,3/3.5/-0.55,4/5.25/0.05}{
  \draw[black!50,decorate,decoration={coil,aspect=0.35,segment length=2.2pt,amplitude=2.2pt}]
    ([yshift=0.75cm]s\i.center) -- ([yshift=0.75cm]s\i.center |- {0,\yy});
  \fill[disk] ([yshift=0.75cm]s\i.center |- {0,\yy}) circle (2.6pt);
}
\node[gridlbl] at (6.3,-6.0) {$\theta{=}k$: in $\{20,40,60,80\}$\\out $\{5,10,200,400\}$ N/m};

% ============ ROW 5: inclined slide ============
\node[rowlbl] at (\rowlabelx,-8.0) {Inclined slide};
\foreach \i/\x in {1/0,2/1.75,3/3.5,4/5.25}{
  \node[framebox] (i\i) at (\x,-8.0) {};
}
\foreach \i in {1,2,3,4}{
  \draw[black!55,line width=1.0pt] ([xshift=-0.7cm,yshift=-0.55cm]i\i.center) -- ([xshift=0.7cm,yshift=0.55cm]i\i.center);
}
\fill[disk] ([xshift=-0.55cm,yshift=-0.43cm]i1.center) circle (2.4pt);
\fill[disk] ([xshift=-0.1cm,yshift=-0.08cm]i2.center) circle (2.4pt);
\fill[disk] ([xshift=0.35cm,yshift=0.27cm]i3.center) circle (2.4pt);
\fill[disk] ([xshift=0.6cm,yshift=0.47cm]i4.center) circle (2.4pt);
\node[gridlbl] at (6.3,-8.0) {$\theta{=}\mu$: in $\{.1,.2,.3,.4\}$\\out $\{.01,.7,1.0\}$ (dimensionless)};

% time axis label under the whole strip
\node[font=\scriptsize,black!45] at (2.6,-9.05) {four representative frames, $t=0 \to$ end of clip};
\draw[-{Stealth[length=1.6mm]},black!35] (-0.75,-8.85) -- (5.9,-8.85);

\end{tikzpicture}
